\documentclass[11pt,a4paper]{article}
\usepackage[margin=18mm,top=17mm,bottom=19mm]{geometry}
\usepackage{amsmath,amssymb,graphicx,array,fancyhdr,lastpage,textcomp}
\usepackage{fontspec}
\setmainfont{texgyreheros-regular.otf}[BoldFont=texgyreheros-bold.otf,ItalicFont=texgyreheros-italic.otf,BoldItalicFont=texgyreheros-bolditalic.otf]
\setlength{\parindent}{0pt}
\setlength{\parskip}{3pt}
\pagestyle{fancy}\fancyhf{}
\renewcommand{\headrulewidth}{0pt}\renewcommand{\footrulewidth}{0.3pt}
\fancyfoot[L]{\footnotesize by paper.shrit.in\quad |\quad normal}
\fancyfoot[R]{\footnotesize Page \thepage\ of \pageref{LastPage}}
\newcommand{\question}[3]{\par\vspace{8pt}\noindent\begin{minipage}[t]{0.075\linewidth}\textbf{#1)}\end{minipage}\begin{minipage}[t]{0.855\linewidth}#3\end{minipage}\hfill\begin{minipage}[t]{0.055\linewidth}\raggedleft(#2)\end{minipage}\par}
\begin{document}
{\LARGE\bfseries Question Paper}\hfill {\small Registration No.: \rule{33mm}{0.3pt}}\par\vspace{8pt}
\begin{center}{\large\bfseries MANIPAL FORMAT | PRACTICE PAPER}\\[6pt]{\bfseries THIRD SEMESTER B.TECH. | MIDSEM PRACTICE}\\[4pt]{\bfseries DIGITAL SYSTEMS AND COMPUTER ORGANIZATION [CSS 2104]}\\[4pt]{\small COMPUTER SCIENCE AND ENGINEERING}\\[3pt]{\small SET B: NEWLY AUTHORED QUESTIONS}\end{center}
\textbf{Marks: 30}\hfill\textbf{Suggested duration: 90 minutes}\par\vspace{3pt}\hrule\vspace{5pt}
{\small\textbf{Answer all questions.} Questions 1--5 are MCQs worth 1 mark each. Select one correct option per MCQ. The remaining questions are theory questions worth 25 marks in total; show working and draw labelled diagrams where appropriate. Modules 1-4; module 5 through memory operations. Independent practice paper; not an official university examination.}\par
\medskip\textbf{Section A: Multiple-choice questions (5 marks)}\par
\question{1}{1}{How many cells are in a four-variable Karnaugh map?\par (A) 4\par (B) 8\par (C) 16\par (D) 32}
\question{2}{1}{How many input variables differ between two adjacent Karnaugh-map cells?\par (A) 1\par (B) 2\par (C) 3\par (D) 4}
\question{3}{1}{Which is a valid group size when minimizing a sum-of-products expression on a Karnaugh map?\par (A) 3\par (B) 5\par (C) 6\par (D) 8}
\question{4}{1}{For $F(A,B,C,D)=\Sigma m(0,2,5,7,8,10,13,15)$, with A most significant, which expression is equivalent?\par (A) $B\oplus D$\par (B) $\neg B\neg D+BD$\par (C) $B+D$\par (D) $BD$}
\question{5}{1}{How can a two-input NAND gate act as an inverter for input X?\par (A) Tie both inputs to 0\par (B) Leave one input open\par (C) Tie both inputs to X\par (D) Tie one input to 0}
\newpage\textbf{Section B: Theory questions (25 marks)}\par
\question{6}{3}{Using 5-bit two\textquotesingle s complement, evaluate (i) $11+7$, (ii) $-9+5$, and (iii) $6-13$. Give the binary result, its signed interpretation and whether signed overflow occurs in each case.}
\question{7}{2}{Implement $F(A,B,C)=A\oplus B\oplus C$ using a 4-to-1 multiplexer. Choose A and B as select inputs and specify $I_0,I_1,I_2,I_3$; an inverter is available.}
\question{8}{5}{Design a synchronous two-bit counter with T flip-flops and enable input E. For $E=1$, the state sequence is $00\to01\to10\to11\to00$; for $E=0$, the state holds. Draw the transition table, derive both T-input equations and draw the circuit.}
\question{9}{3}{Write synthesizable Verilog for a four-bit ring counter with positive-edge clock and active-high synchronous reset. Reset must load \texttt{0001}; subsequent clocks rotate the one toward the next more significant bit, wrapping from bit 3 to bit 0.}
\question{10}{2}{Explain why an asynchronous ripple counter can briefly produce incorrect decoded outputs. Contrast its clock connection with that of a synchronous counter.}
\question{11}{5}{Design a two-bit unsigned magnitude comparator for $A=a_1a_0$ and $B=b_1b_0$. Derive Boolean expressions for $A>B$, $A=B$ and $A<B$. Draw a gate-level diagram and verify the outputs for $A=10_2$, $B=01_2$.}
\question{12}{3}{A byte-addressable memory stores 32-bit words. How many words fit in 4 KiB? How many address bits select individual bytes? If a word starts at hexadecimal address \texttt{0x1000}, list the byte addresses it occupies and the starting address of the next word.}
\question{13}{2}{Describe a memory-read operation using the memory address register (MAR), memory data register (MDR) and read control signal. State what changes for a memory-write operation.}
\vfill\begin{center}\small End of question paper\end{center}\end{document}